% TOP_PROBLEM: {"id": 2800701, "title": "10 Lectures and 42 Open Problems — Gilbert-Varshamov bound", "category_id": 15, "category": {"id": 15, "name": "computer_science", "display_name": "Computer Science", "description": "Computational complexity, algorithms, and theoretical CS.", "slug": "computer-science", "order_index": 15, "created_at": "2026-07-31T15:26:25.671Z"}, "status": "partially_solved", "problem_number": "AMR-027-0701", "set_id": 13}
% FIRST_POSTED: 2026-10-02
% ENTRY_KIND: statement_only
% UnsolvedMath ID 2800701; source catalogue code AMR-027-0701
% !TeX program = lualatex
% Definition and sources only; this file is not an LLM solution attempt.
\documentclass[11pt,a4paper]{article}
\usepackage[margin=25mm]{geometry}
\usepackage{fontspec}
\setmainfont{lmroman10-regular.otf}[BoldFont=lmroman10-bold.otf,
 ItalicFont=lmroman10-italic.otf,BoldItalicFont=lmroman10-bolditalic.otf]
\usepackage{amsmath,amssymb,mathtools}
\usepackage{unicode-math}
\setmathfont{latinmodern-math.otf}
\AtBeginDocument{\let\setminus\smallsetminus}
\usepackage{xurl,hyperref}
\hypersetup{colorlinks=true,urlcolor=blue}
\setcounter{secnumdepth}{0}
\setlength{\parindent}{0pt}
\setlength{\parskip}{5pt}
\setlength{\emergencystretch}{3em}
\begin{document}
\section*{10 Lectures and 42 Open Problems — Gilbert-Varshamov bound}\label{2800701}
{\small UnsolvedMath ID 2800701 · AMR-027-0701 · Computer Science · AMR Open Problem Lists}
\subsection{Definitions and mathematical statement}
A binary code of block length $n$ is a subset $C\subseteq\{0,1\}^n$.
Its minimum distance is the least Hamming distance between distinct
codewords, where Hamming distance counts differing coordinates.
Write $A_2(n,d)$ for the largest size of a binary code with minimum
distance at least $d$, and define
\[
 R_2(\delta)=\limsup_{n\to\infty}
            \frac1n\log_2 A_2(n,\lceil\delta n\rceil),
 \qquad 0<\delta<\tfrac12.
\]
The binary entropy function is
$H_2(\delta)=-\delta\log_2\delta-(1-\delta)\log_2(1-\delta)$.
The Gilbert--Varshamov bound asserts $R_2(\delta)\ge1-H_2(\delta)$.

The source asks two distinct questions. First, construct deterministic
explicit code families attaining this rate: for each fixed $\delta$
and $\varepsilon>0$, obtain arbitrarily large block lengths with
relative distance at least $\delta$ and rate at least
$1-H_2(\delta)-\varepsilon$, together with a construction and encoding
procedure polynomial in the block length. A linear construction can
be specified by a generator matrix; nonlinear codes require another
succinct, efficiently evaluable encoder.

Second, is $R_2(\delta)=1-H_2(\delta)$ throughout this interval, or can
binary codes asymptotically exceed the bound? This extremal question
allows nonlinear codes and places no computational restriction on
them. Efficient decoding is a separate desirable property, not an
extra hypothesis of the two stated questions. Exponential enumeration
of all codewords does not supply the requested efficient explicitness.
\subsection{Short English statement}
Construct explicit binary codes attaining the Gilbert–Varshamov rate, and determine whether that rate is optimal among all binary codes.
\subsection{Sources}
\begin{itemize}
\item[S1] ULAM.AI and the UnsolvedMath Contributors, supplied problems.json, record 2800701 (AMR-027-0701), AMR Open Problem Lists. Catalogue identity and source transcription; CC BY 4.0.
\url{https://huggingface.co/datasets/ulamai/UnsolvedMath}.
\item[S2] Bandeira - 42 Open Problems in Data Science (2016), Open Problem 7.1. Original source indexed by AMR-027-0701.
\url{https://afonsobandeira.wordpress.com/2015/11/27/10l42grouptesting/}.
\item[S3] A. S. Bandeira, Ten Lectures and Forty-Two Open Problems in the Mathematics of Data Science, corresponding numbered problem and surrounding definitions.
\url{https://people.math.ethz.ch/~abandeira/TenLecturesFortyTwoProblems.pdf}.
\end{itemize}
\end{document}
